\section{Related Literature}
\label{sec:literature}

This paper sits at the intersection of four literatures.

\paragraph{Bid-coordination tests and bidder classification.}
\citet{bajari2003deciding} develop powerful tests for
collusion---exchangeability of bid functions and conditional
independence of bid residuals---but require an \emph{ex ante}
partition of bidders into ``suspected colluders'' and a ``competitive
fringe.'' No existing method provides this partition from data alone.
Subsequent work detects bid rotation
\citep{porter1993detection, porter1999ohio}, collusion in timber
auctions \citep{baldwin1997bidder}, bidder groups from co-bidding
patterns \citep{conley2016detecting}, and collusion dynamics in
long-lived cartels \citep{clark2021collusion,
schurter2020identification}. \citet{conley2016detecting} are the
closest precedent in data requirements, though they detect group
structure rather than flag individual firms. In every case the
partition is constructed \emph{ex post} or by assumption. The FL
classification provides a candidate for this partition: a data-driven,
\emph{ex ante} classification grounded in a participation anomaly
that can feed into Bajari--Ye tests (Section~\ref{sec:results_bajari}).

\paragraph{Proactive screens for bid rigging.}
\citet{imhof2019detecting} and \citet{imhof2018screening} train ML
classifiers on tender-level features and achieve high accuracy against
known cartels; see also \citet{abrantesmetz2006a},
\citet{harrington2008detecting}, \citet{chassang2022robust}, and
\citet{wallimann2023machine}. All share two requirements: granular bid
values and tender-level units of analysis. Many procurement systems,
particularly in developing economies, record who bid and who won but
not the bid amounts. The FL screen operates at the \emph{firm} level
using only participation and outcome records, enabling a two-stage
triage workflow: FL flags prioritize environments for closer scrutiny,
and bid-level analysis follows on the flagged subset. In a direct
comparison on the same ground truth, FL achieves AUC~$= 0.94$ versus
0.79 for an Imhof-style CV proxy, and the two screens are nearly
orthogonal (correlation 0.06).%
\footnote{Our Imhof comparison uses a CV median split, a coarser proxy
for the full implementation; see Section~\ref{sec:results_detection}.}

\paragraph{Cover bidding.}
Cover bidders have been documented in prosecuted cartels
\citep{porter1999ohio, pesendorfer2000study, asker2010leniency,
marshall2012economics} and in theoretical comparisons of auction
formats \citep{athey2011comparing}. In every case, identification is
\emph{ex post}---the cover bidder is named after the cartel is
caught. Our conceptual framework models cover bidding \emph{ex ante},
distinguishing two regimes with different bid distributions and
different detection vulnerabilities. Under the coordinated regime,
dispersion-based screens lose power, highlighting a setting in which a
participation-based screen may be especially useful.

\paragraph{Procurement in developing economies.}
The bid-rigging literature concentrates on developed economies, yet
the stakes are highest where enforcement resources are scarce and
investigative capacity is limited
\citep{sanchezgraells2019screening}. The FL screen's minimal data
requirements---participation records and outcomes, without bid
microdata---make it directly portable to low-capacity procurement
systems. BEC's scale (4.5 million tender-items, 40 million bids,
11 years, two auction formats) provides a demanding test bed for a
screen designed to work with less.
